\chapter{可信性能测量基础}

\section{本章要解决的问题}

从这一章开始，本书进入性能工程部分。很多人以为性能优化的第一步是“改代码”，但真正可靠的第一步是“知道自己测到的是什么”。如果测量方法不可信，后面所有结论都可能是错的：你以为某个优化快了 30\%，实际可能是编译器删掉了循环；你以为某个版本慢了，实际可能是 CPU 降频、cache 状态、OS 调度、页错误、随机输入或 I/O 噪声在影响。

本章要建立一套性能测量的底层心智模型。你需要从一开始就把 benchmark 看成一次受控实验，而不是一次随手计时。

一个可信 benchmark 至少要回答：

\begin{itemize}
  \item 被测对象是什么？是一个函数、一个循环、一个完整程序，还是一次系统调用？
  \item 输入是什么？规模、分布、对齐、cache 状态、随机种子是否明确？
  \item 计时区域在哪里？是否混入了初始化、分配、I/O、随机数生成、校验、日志输出？
  \item 编译器是否可能把被测代码删除、合并、常量折叠或移动？
  \item 结果是否先证明正确，再讨论性能？
  \item 测量重复了多少次？报告的是最小值、均值、中位数还是分位数？
  \item 环境是否记录？CPU 型号、频率策略、编译器版本、编译参数、操作系统状态是否可复现？
\end{itemize}

本章会先讲最核心的机制，后续第 16 章再系统讲 benchmark 设计，第 17 章再深入统计与噪声分析。你现在要掌握的是：看到一个性能数字时，能判断它是否有资格作为工程证据。

\begin{keyidea}
性能测量不是“得到一个时间数字”，而是建立一条可复核的证据链：正确性、被测工作量、时间源、控制变量、重复测量、统计口径、二进制审计和环境记录。
\end{keyidea}

\section{为什么计时不等于 benchmark}

先看一个常见错误：

\begin{lstlisting}[language=C++]
#include <chrono>
#include <cstdint>
#include <cstdio>

int main()
{
    auto begin = std::chrono::high_resolution_clock::now();

    std::uint64_t sum = 0;
    for (std::uint64_t i = 0; i < 1000000000ULL; ++i) {
        sum += i;
    }

    auto end = std::chrono::high_resolution_clock::now();
    auto ns = std::chrono::duration_cast<std::chrono::nanoseconds>(end - begin);
    std::printf("%lld\n", static_cast<long long>(ns.count()));
}
\end{lstlisting}

这个程序看起来在测 10 亿次加法。实际上它可能什么都没测到。优化器可以推导：

\begin{lstlisting}
sum = 0 + 1 + 2 + ... + 999999999
\end{lstlisting}

如果 \code{sum} 没有被观察，甚至整个循环都可能被删除。就算打印了 \code{sum}，编译器也可能用公式替代循环。因为从 C++ 抽象机视角看，只要可观察行为相同，编译器可以自由改写。

另一个错误：

\begin{lstlisting}[language=C++]
auto begin = clock_now();
std::vector<float> x(n);
fill_random(x);
auto result = kernel(x);
std::printf("%f\n", result);
auto end = clock_now();
\end{lstlisting}

这里计时区域包含了：

\begin{lstlisting}
allocation + initialization + random generation + kernel + printf
\end{lstlisting}

如果目标是测 \code{kernel}，这个数字没有表达清楚。它可能主要测到了内存分配器、随机数生成器、首次触页、I/O flush 或动态链接的 lazy binding。

\begin{warningbox}
“程序运行了多久”和“kernel 花了多久”不是同一个问题。高性能计算和 AI 算子优化通常关心热路径 kernel 的稳定吞吐；计时区域必须把 setup、timed region、validation 分开。
\end{warningbox}

\section{性能实验的基本结构}

一个最小可信实验可以拆成五段：

\begin{lstlisting}
1. describe environment:
       CPU, compiler, flags, OS, governor, command

2. setup:
       allocate data, initialize input, pin thread if needed

3. warmup:
       run kernel without recording final timing

4. timed region:
       measure only the intended workload

5. validation and report:
       compare with reference, compute statistics, save raw data
\end{lstlisting}

注意顺序：正确性在性能之前。如果一个优化版本结果错了，它再快也没有意义。

\begin{mentalmodel}
benchmark 是一个实验系统。kernel 只是其中一个部件；输入、环境、计时器、统计、报告和二进制审计都是实验系统的一部分。训练自己时，不要只问“快不快”，还要问“我凭什么相信这个数字”。
\end{mentalmodel}

\section{被测对象：workload、kernel、measurement}

先定义几个词。

\begin{center}
\begin{tabular}{p{0.30\linewidth}p{0.54\linewidth}}
\toprule
术语 & 本书中的含义 \\
\midrule
\term{workload}{工作负载} & 数据规模、数据分布、访问模式和运行条件。 \\
\term{kernel}{核心计算} & 希望重点优化和计时的函数或循环。 \\
\term{timed region}{计时区域} & 实际包在计时器之间的代码范围。 \\
\term{harness}{测试框架} & 负责 setup、warmup、重复测量、统计和报告。 \\
\term{reference}{参考实现} & 语义清晰、容易验证正确性的版本，用来检查优化版本。 \\
\bottomrule
\end{tabular}
\end{center}

例子：

\begin{lstlisting}[language=C++]
float dot_product_ref(std::span<const float> a, std::span<const float> b)
{
    float sum = 0.0f;
    for (std::size_t i = 0; i < a.size(); ++i) {
        sum += a[i] * b[i];
    }
    return sum;
}
\end{lstlisting}

这里 \code{dot_product_ref} 是 reference 或 baseline，不一定是最终 kernel。workload 包括：

\begin{itemize}
  \item \code{a} 和 \code{b} 的长度。
  \item 数据是否随机、全零、全一、递增、含 NaN、含 denormal。
  \item 指针是否对齐。
  \item 数据是否已经在 cache 中。
  \item 是否只测一个长度，还是扫描多个长度。
\end{itemize}

如果报告只写“dot product 需要 3 ms”，信息不够。至少应写：

\begin{lstlisting}
kernel:
    dot_product_ref(float*, float*, n)

workload:
    n = 1048576
    data = uniform random in [-1, 1]
    arrays allocated once, reused
    warm-cache repeated measurement

compiler:
    clang++ -O3 -march=native -DNDEBUG

machine:
    CPU model, core count, governor, OS kernel

metric:
    median nanoseconds over 31 measured repetitions
\end{lstlisting}

\section{正确性先于性能}

性能优化常见灾难是“优化版本快了，因为少算了”。因此每个 benchmark 都要先有 correctness reference。

以整数求和为例：

\begin{lstlisting}[language=C++]
#include <cstdint>
#include <span>

std::uint64_t sum_u64_ref(std::span<const std::uint64_t> values)
{
    std::uint64_t sum = 0;
    for (const std::uint64_t value : values) {
        sum += value;
    }
    return sum;
}
\end{lstlisting}

为什么用 \code{std::uint64_t}？因为 unsigned overflow 在 C++ 中定义为模 \code{2^64} 回绕。若使用 \code{std::int64_t}，输入过大时 signed overflow 是未定义行为，reference 也可能被优化器按“不会溢出”的假设改写。

测试思路：

\begin{lstlisting}[language=C++]
void check_sum_kernel(std::span<const std::uint64_t> values)
{
    const std::uint64_t expected = sum_u64_ref(values);
    const std::uint64_t actual = sum_u64_optimized(values);
    if (actual != expected) {
        std::abort();
    }
}
\end{lstlisting}

如果是浮点，正确性不能总用完全相等。应根据算法和数值误差选择容忍度：

\begin{lstlisting}
absolute error:
    abs(actual - expected) <= epsilon

relative error:
    abs(actual - expected) <= epsilon * max(abs(expected), 1)
\end{lstlisting}

\begin{deepdive}
浮点优化不仅是性能问题，也是语义问题。改变求和顺序会改变舍入误差。

启用 \code{-ffast-math} 还会改变 NaN、Inf、signed zero 和结合律相关的语义边界。后续讲 fast-math 和 AI 算子时会专门展开。现在先记住：浮点 benchmark 必须同时报告数值误差策略。
\end{deepdive}

\section{时间源：你到底在读什么钟}

性能测量常见时间源有几类。

\begin{center}
\begin{tabular}{p{0.25\linewidth}p{0.61\linewidth}}
\toprule
时间源 & 适合用途和主要风险 \\
\midrule
\code{steady_clock} & C++ 通用单调墙钟计时；分辨率和开销依平台而异。 \\
\code{clock_gettime} & Linux 下可选择 \code{CLOCK_MONOTONIC} 等时间源；系统调用或 vDSO 路径需要了解。 \\
TSC & 低开销 cycle-like 计数；风险是乱序、迁核、频率解释和序列化问题。 \\
PMU counters & 观察硬件事件，例如 cycles、instructions、cache miss；风险是权限、事件语义、复用和平台差异。 \\
\code{/usr/bin/time} & 粗粒度进程级时间和资源统计；不适合细粒度 kernel。 \\
\bottomrule
\end{tabular}
\end{center}

\subsection{\code{steady_clock}}

C++ 中应优先使用 \code{std::chrono::steady_clock} 作为入门计时器，因为它承诺单调不倒退。

\begin{lstlisting}[language=C++]
using Clock = std::chrono::steady_clock;

const auto begin = Clock::now();
run_kernel();
const auto end = Clock::now();

const auto elapsed_ns =
    std::chrono::duration_cast<std::chrono::nanoseconds>(end - begin).count();
\end{lstlisting}

不要默认 \code{high_resolution_clock} 一定更好。标准没有保证它一定是单调钟；在某些实现中它只是其他 clock 的别名。入门阶段用 \code{steady_clock} 更稳。

\subsection{TSC 与 cycle 的关系}

x86 上可以用 \instruction{rdtsc} 或 \instruction{rdtscp} 读取 time-stamp counter。它返回一个 64 位计数，常见写法是 \register{edx}:\register{eax} 组合：

\begin{lstlisting}[language=C++]
#include <cstdint>
#include <x86intrin.h>

std::uint64_t read_tsc()
{
    return __rdtsc();
}
\end{lstlisting}

但是 TSC 不是新手测量的万能工具。风险包括：

\begin{itemize}
  \item \instruction{rdtsc} 不是完全序列化指令，周围指令可能乱序跨过它。
  \item 如果线程迁移到另一个 core，早期或特殊系统上的 TSC 同步可能成为问题。
  \item 现代 CPU 的 TSC 常以 invariant rate 增长，不一定等于当前核心实际执行频率。
  \item Turbo、降频、C-state 会让“每秒多少 TSC tick”和“每条指令需要多少 core cycle”的解释变复杂。
\end{itemize}

更严谨的 TSC 计时需要 fence：

\begin{lstlisting}[language=C++]
#include <cstdint>
#include <x86intrin.h>

std::uint64_t read_tsc_begin()
{
    _mm_lfence();
    const std::uint64_t value = __rdtsc();
    _mm_lfence();
    return value;
}

std::uint64_t read_tsc_end()
{
    _mm_lfence();
    const std::uint64_t value = __rdtsc();
    _mm_lfence();
    return value;
}
\end{lstlisting}

不同平台和论文中还会看到 \instruction{cpuid}、\instruction{rdtscp}、\instruction{lfence} 的不同组合。本书训练阶段的原则是：

\begin{itemize}
  \item 初期先用 \code{steady_clock} 建立正确实验结构。
  \item 需要 cycle 级分析时，再引入 TSC，并明确 fence 策略。
  \item 真正分析微架构时，优先结合 PMU counter，而不是只看 TSC。
\end{itemize}

\begin{warningbox}
不要把 TSC delta 直接称为“CPU cycles”而不解释。很多现代机器上 TSC 是稳定参考计数，核心实际频率会变化。报告时可以写 \code{TSC ticks} 或说明换算假设。
\end{warningbox}

\section{编译器会删除你的 benchmark}

优化器只需要保留 C++ 抽象机的可观察行为。benchmark 里最常见的错误是被测计算没有可观察结果。

\subsection{dead-code elimination}

错误：

\begin{lstlisting}[language=C++]
void benchmark_bad(std::span<const std::uint64_t> values)
{
    std::uint64_t sum = 0;
    for (const std::uint64_t value : values) {
        sum += value;
    }
}
\end{lstlisting}

\code{sum} 没有被使用。优化器可以删除整个循环。

最简单的防护是把结果返回给调用者：

\begin{lstlisting}[language=C++]
std::uint64_t benchmark_better(std::span<const std::uint64_t> values)
{
    std::uint64_t sum = 0;
    for (const std::uint64_t value : values) {
        sum += value;
    }
    return sum;
}
\end{lstlisting}

但如果调用者也没有使用返回值，仍可能被删除。benchmark harness 需要显式制造可观察使用。

\subsection{\code{do_not_optimize}}

GCC/Clang 常用一个空 inline asm 阻止优化器认为值无用：

\begin{lstlisting}[language=C++]
template <class T>
inline void do_not_optimize(const T& value)
{
    asm volatile("" : : "g"(value) : "memory");
}
\end{lstlisting}

用法：

\begin{lstlisting}[language=C++]
const std::uint64_t result = benchmark_better(values);
do_not_optimize(result);
\end{lstlisting}

这段 inline asm 没有机器指令，但它告诉编译器：\code{value} 被某个不可见的 asm 使用，不能随便删除产生它的计算。\code{"memory"} 还限制周围内存访问重排。

\begin{warningbox}
\code{do_not_optimize} 是 benchmark 工具，不是业务代码工具。它的目的是约束编译器优化，让实验可测；不要把它放进生产热路径。
\end{warningbox}

\subsection{\code{clobber_memory}}

另一个常用工具：

\begin{lstlisting}[language=C++]
inline void clobber_memory()
{
    asm volatile("" : : : "memory");
}
\end{lstlisting}

它告诉编译器：内存状态可能被未知 asm 改变，因此不要把内存读写随便跨过这个点移动。它不会清空 CPU cache，也不是硬件内存栅栏。

\subsection{constant folding}

错误：

\begin{lstlisting}[language=C++]
std::uint64_t x = expensive(123);
\end{lstlisting}

如果 \code{expensive(123)} 在编译期可推导，优化器可能直接替换成常量。benchmark 输入应来自运行时数据，或者至少来自编译器无法完全证明的内存。

例子：

\begin{lstlisting}[language=C++]
std::vector<std::uint64_t> make_input(std::size_t size)
{
    constexpr std::uint64_t INPUT_MULTIPLIER = 6364136223846793005ULL;
    constexpr std::uint64_t INPUT_INCREMENT = 1442695040888963407ULL;

    std::vector<std::uint64_t> values(size);
    std::uint64_t state = 1;
    for (std::uint64_t& value : values) {
        state = state * INPUT_MULTIPLIER + INPUT_INCREMENT;
        value = state;
    }
    return values;
}
\end{lstlisting}

这里用简单线性同余生成输入。它不是高质量随机数生成器，但适合构造可复现、运行时生成的数据。真正 benchmark 时，输入生成仍应放在 setup，不放入 timed region。

\section{计时区域必须隔离}

错误结构：

\begin{lstlisting}[language=C++]
auto begin = Clock::now();
std::vector<std::uint64_t> values = make_input(size);
const std::uint64_t result = sum_u64(values);
check(result);
auto end = Clock::now();
\end{lstlisting}

这测的是：

\begin{lstlisting}
allocation + input generation + sum + check
\end{lstlisting}

正确结构：

\begin{lstlisting}[language=C++]
std::vector<std::uint64_t> values = make_input(size);
const std::uint64_t expected = sum_u64_ref(values);

auto begin = Clock::now();
const std::uint64_t result = sum_u64(values);
auto end = Clock::now();

if (result != expected) {
    std::abort();
}
\end{lstlisting}

如果要重复测量：

\begin{lstlisting}
setup:
    allocate and initialize input once
    compute expected reference once

warmup:
    run kernel several times

measurement:
    for each repetition:
        begin timer
        run kernel
        end timer
        record elapsed
        consume result

validation:
    compare result with expected
\end{lstlisting}

\begin{deepdive}
有些 workload 本身包含分配或 I/O，比如数据库查询、HTTP 服务、编译器编译一个文件。那不是不能测，而是要诚实命名 benchmark：你测的是 end-to-end workload，不是纯计算 kernel。不同问题需要不同边界。
\end{deepdive}

\section{warmup：第一次运行不代表稳态}

第一次运行常常慢，因为它可能包含：

\begin{itemize}
  \item 代码页和数据页首次触页。
  \item instruction cache 和 data cache 冷启动。
  \item branch predictor 尚未学习输入模式。
  \item 动态库 lazy binding。
  \item CPU 从低功耗状态升频。
  \item 内存分配器初始化。
\end{itemize}

所以 benchmark 通常先 warmup：

\begin{lstlisting}[language=C++]
constexpr int BENCHMARK_WARMUP_REPETITIONS = 5;

for (int iter = 0; iter < BENCHMARK_WARMUP_REPETITIONS; ++iter) {
    const std::uint64_t result = sum_u64(values);
    do_not_optimize(result);
}
\end{lstlisting}

warmup 不是为了“作弊让 cache 热起来”，而是为了明确你测的是某个状态。你可以测 cold-cache，也可以测 warm-cache，但必须说清楚。

\begin{center}
\begin{tabular}{p{0.24\linewidth}p{0.60\linewidth}}
\toprule
状态 & 含义 \\
\midrule
cold-cache & 每次计时前尽量让数据不在 cache 中，模拟首次访问或大工作集。 \\
warm-cache & 同一数据重复运行，测稳定热路径吞吐。 \\
steady-state & 程序已经过初始化、预测器和 cache 已接近稳定状态。 \\
\bottomrule
\end{tabular}
\end{center}

\section{重复测量和统计口径}

单次测量几乎没有说服力。需要多次运行并保存原始数据。

常见统计量：

\begin{center}
\begin{tabular}{p{0.16\linewidth}p{0.38\linewidth}p{0.30\linewidth}}
\toprule
统计量 & 含义 & 风险 \\
\midrule
min & 最小观测值，接近“噪声最少时”的能力。 & 可能过于乐观。 \\
mean & 平均值。 & 对 outlier 敏感。 \\
median & 中位数，代表典型运行。 & 不展示尾部风险。 \\
p95 & 95 分位，展示较慢尾部。 & 需要足够样本。 \\
max & 最慢观测。 & 可能只是一次系统干扰。 \\
\bottomrule
\end{tabular}
\end{center}

HPC kernel 优化常看 min 或 median，因为目标是估计稳定热路径上限。服务端延迟常看 p95、p99，因为用户体验受尾延迟影响。AI 算子库通常会同时报告 median、min 和标准差或变异系数。

\begin{warningbox}
不要只报告一个数字且不说明它是什么。写“用时 1.2 ms”是不完整的；应写“31 次重复，median = 1.2 ms，min = 1.17 ms，p95 = 1.35 ms”。
\end{warningbox}

\section{一个最小 C++20 benchmark harness}

下面给一个教材用最小 harness。它不是要替代 Google Benchmark，而是让你看清楚每个部件为什么存在。

\begin{lstlisting}[language=C++]
#include <algorithm>
#include <chrono>
#include <cstdint>
#include <cstdio>
#include <cstdlib>
#include <span>
#include <vector>

namespace {

constexpr std::size_t BENCHMARK_INPUT_SIZE = 1u << 20;
constexpr int BENCHMARK_WARMUP_REPETITIONS = 5;
constexpr int BENCHMARK_MEASURED_REPETITIONS = 31;
constexpr std::uint64_t INPUT_MULTIPLIER = 6364136223846793005ULL;
constexpr std::uint64_t INPUT_INCREMENT = 1442695040888963407ULL;

using Clock = std::chrono::steady_clock;

template <class T>
inline void do_not_optimize(const T& value)
{
    asm volatile("" : : "g"(value) : "memory");
}

std::vector<std::uint64_t> make_input(std::size_t size)
{
    std::vector<std::uint64_t> values(size);
    std::uint64_t state = 1;
    for (std::uint64_t& value : values) {
        state = state * INPUT_MULTIPLIER + INPUT_INCREMENT;
        value = state;
    }
    return values;
}

std::uint64_t sum_u64(std::span<const std::uint64_t> values)
{
    std::uint64_t sum = 0;
    for (const std::uint64_t value : values) {
        sum += value;
    }
    return sum;
}

std::uint64_t elapsed_ns_for_one_run(std::span<const std::uint64_t> values,
                                     std::uint64_t* result_out)
{
    const auto begin = Clock::now();
    const std::uint64_t result = sum_u64(values);
    const auto end = Clock::now();

    *result_out = result;
    do_not_optimize(result);

    return static_cast<std::uint64_t>(
        std::chrono::duration_cast<std::chrono::nanoseconds>(end - begin).count());
}

double median_ns(std::vector<std::uint64_t> samples)
{
    std::sort(samples.begin(), samples.end());
    const std::size_t middle = samples.size() / 2;
    return static_cast<double>(samples[middle]);
}

}  // namespace

int main()
{
    const std::vector<std::uint64_t> values = make_input(BENCHMARK_INPUT_SIZE);
    const std::uint64_t expected = sum_u64(values);

    for (int iter = 0; iter < BENCHMARK_WARMUP_REPETITIONS; ++iter) {
        const std::uint64_t result = sum_u64(values);
        do_not_optimize(result);
    }

    std::vector<std::uint64_t> samples;
    samples.reserve(BENCHMARK_MEASURED_REPETITIONS);

    for (int iter = 0; iter < BENCHMARK_MEASURED_REPETITIONS; ++iter) {
        std::uint64_t result = 0;
        const std::uint64_t elapsed_ns = elapsed_ns_for_one_run(values, &result);
        if (result != expected) {
            std::abort();
        }
        samples.push_back(elapsed_ns);
    }

    const auto [min_it, max_it] = std::minmax_element(samples.begin(), samples.end());
    std::printf("n=%zu min_ns=%llu median_ns=%.0f max_ns=%llu\n",
                values.size(),
                static_cast<unsigned long long>(*min_it),
                median_ns(samples),
                static_cast<unsigned long long>(*max_it));
}
\end{lstlisting}

编译：

\begin{lstlisting}[language=bash]
clang++ -std=c++20 -O3 -march=native -DNDEBUG \
    bench_sum.cpp -o bench_sum
./bench_sum
\end{lstlisting}

这个 harness 做对了几件事：

\begin{itemize}
  \item 输入在计时前生成。
  \item reference 结果在计时前计算。
  \item warmup 和 measurement 分离。
  \item 每次计时只包住 \code{sum_u64}。
  \item 结果被校验并被 \code{do_not_optimize} 消费。
  \item 重复测量并报告 min、median、max。
\end{itemize}

它还不够完整：

\begin{itemize}
  \item 没有输出原始 CSV。
  \item 没有记录环境。
  \item 没有 CPU affinity。
  \item 没有 PMU counter。
  \item 没有多输入规模扫描。
  \item 没有统计 p95 和变异系数。
\end{itemize}

这些会在后续章节逐步补齐。

\section{环境控制：CPU 不只是在跑你的代码}

一次 benchmark 运行时，机器上还有很多因素：

\begin{itemize}
  \item 操作系统可能调度其他进程。
  \item 中断可能打断当前 core。
  \item 线程可能迁移到另一个 core。
  \item CPU 可能升频或降频。
  \item SMT sibling 可能共享执行资源。
  \item 后台服务可能污染 cache 或占用内存带宽。
  \item 虚拟机或容器可能隐藏真实硬件行为。
\end{itemize}

入门阶段可以先做低成本控制：

\begin{lstlisting}[language=bash]
# 查看 CPU 型号
lscpu

# 查看当前频率策略，具体路径依发行版而异
cat /sys/devices/system/cpu/cpu0/cpufreq/scaling_governor

# 绑定到某个 CPU 运行
taskset -c 2 ./bench_sum

# 粗略查看进程级时间
/usr/bin/time -v ./bench_sum
\end{lstlisting}

报告中至少记录：

\begin{lstlisting}
CPU:
    model name, cores, SMT on/off if known

OS:
    kernel version, distribution

compiler:
    compiler name and version
    flags

run command:
    exact command line

environment:
    taskset core if used
    governor if known
    whether running on laptop battery, VM, container, remote server
\end{lstlisting}

\begin{deepdive}
严格控制频率和隔离 CPU core 可能需要 root 权限和系统配置，例如 performance governor、isolcpus、关闭 turbo、关闭 SMT、固定 NUMA 节点。训练早期不一定要全部做到，但必须知道没有控制意味着结果会有更大噪声。
\end{deepdive}

\section{cache 状态会改变结论}

同一个内存 kernel，在不同工作集大小下可能完全不同。

例子：

\begin{lstlisting}
sum array of uint64:
    n = 1024          likely L1/L2 resident
    n = 1048576       likely LLC or memory bandwidth
    n = 134217728     definitely memory streaming and paging pressure
\end{lstlisting}

如果只测一个 \code{n}，你可能误以为某个优化“总是快”。实际上：

\begin{itemize}
  \item 小数组可能受函数调用、循环开销、前端和分支影响。
  \item 中等数组可能受 cache bandwidth 和 latency 影响。
  \item 大数组可能受内存带宽、TLB、预取器和 NUMA 影响。
\end{itemize}

所以性能报告应把输入规模当作一维变量，而不是固定一个点：

\begin{lstlisting}
n, median_ns, bytes, bandwidth_GBps
1024, ...
4096, ...
16384, ...
65536, ...
1048576, ...
16777216, ...
\end{lstlisting}

带宽计算：

\begin{lstlisting}
bytes_read = n * sizeof(uint64_t)
bandwidth_GBps = bytes_read / seconds / 1e9
\end{lstlisting}

注意这个公式只计算源码层显式读取字节数。真实硬件可能按 cache line 传输，可能有 write allocate，可能有预取，可能有额外 page walk。后面内存层次章节会更精确。

\section{PMU：从时间走向硬件证据}

时间告诉你“慢了多少”，PMU counter 帮你猜“为什么慢”。

Linux 下可以先用：

\begin{lstlisting}[language=bash]
perf stat ./bench_sum
\end{lstlisting}

更具体：

\begin{lstlisting}[language=bash]
perf stat -r 10 \
    -e cycles,instructions,branches,branch-misses,cache-references,cache-misses \
    ./bench_sum
\end{lstlisting}

常见指标：

\begin{center}
\begin{tabular}{p{0.22\linewidth}p{0.31\linewidth}p{0.31\linewidth}}
\toprule
指标 & 直觉 & 注意 \\
\midrule
\code{cycles} & 花了多少周期。 & 受频率、停顿和事件定义影响。 \\
\code{instructions} & retired instructions 数。 & 指令少不一定快。 \\
\code{IPC} & instructions / cycles。 & 高 IPC 不一定代表高性能。 \\
\code{branch-misses} & 分支预测失败。 & 和输入分布强相关。 \\
\code{cache-misses} & cache miss 事件。 & 不同 CPU 事件语义不同。 \\
\bottomrule
\end{tabular}
\end{center}

如果权限不足，你可能看到：

\begin{lstlisting}
No permission to enable cycles event.
\end{lstlisting}

这不是 benchmark 失败，而是环境限制。报告里应写清楚，并改用可用工具：

\begin{lstlisting}[language=bash]
cat /proc/sys/kernel/perf_event_paranoid
/usr/bin/time -v ./bench_sum
\end{lstlisting}

\begin{warningbox}
PMU counter 不是绝对真理。事件可能被复用，事件名称和语义依 CPU 而异，虚拟机里可能不准确。PMU 的正确用法是和源码、汇编、输入规模、时间数据一起形成一致解释。
\end{warningbox}

\section{从源码到汇编再到测量}

性能测量不能只看源码。至少要审计核心循环是否存在。

命令：

\begin{lstlisting}[language=bash]
clang++ -std=c++20 -O3 -march=native -S -masm=intel bench_sum.cpp -o bench_sum.s
clang++ -std=c++20 -O3 -march=native bench_sum.cpp -o bench_sum
objdump -drwC -Mintel bench_sum > bench_sum.objdump.txt
rg -n "sum_u64|add|vpadd|ret" bench_sum.objdump.txt
\end{lstlisting}

你要看：

\begin{itemize}
  \item \code{sum_u64} 是否被内联。
  \item 循环是否还存在。
  \item 是否被向量化。
  \item 是否出现预期加载。
  \item timed region 是否包含你以为的函数。
\end{itemize}

如果函数被内联，不一定是坏事。真实 \code{-O3} 性能经常依赖内联。但报告必须说明你测的是内联后的热路径，而不是一个独立函数调用。

\begin{deepdive}
之后讲编译器优化时，你会看到 scalar replacement、loop unrolling、vectorization、LICM、DCE、constant propagation 等优化。benchmark 的可怕之处在于：这些优化可能正是你想研究的对象，也可能是把实验悄悄变成另一个实验的干扰因素。
\end{deepdive}

\section{性能数字的单位}

不同 kernel 应选择不同单位。

\begin{center}
\begin{tabular}{p{0.20\linewidth}p{0.34\linewidth}p{0.30\linewidth}}
\toprule
单位 & 适用场景 & 例子 \\
\midrule
ns/op & 小操作、函数调用、短 kernel。 & 每次 hash、每次分配。 \\
cycles/op & 微架构分析。 & 每元素周期数。 \\
GB/s & 内存带宽型 kernel。 & copy、fill、sum、transpose。 \\
GFLOP/s & 浮点计算型 kernel。 & dot、GEMM、卷积。 \\
items/s & 解析、压缩、推理 token。 & rows/s、tokens/s。 \\
\bottomrule
\end{tabular}
\end{center}

例子：sum \code{uint64_t}：

\begin{lstlisting}
elements = n
bytes_read = n * 8
seconds = median_ns / 1e9
GB/s = bytes_read / seconds / 1e9
ns/element = median_ns / n
\end{lstlisting}

例子：dot product：

\begin{lstlisting}
for each element:
    one multiply
    one add

FLOPs = 2 * n
GFLOP/s = FLOPs / seconds / 1e9
\end{lstlisting}

要小心 FLOP 计数。FMA 指令 \instruction{vfmadd} 通常按 2 FLOPs 计，因为它完成乘加两个浮点操作。整数算子、量化算子、softmax、layernorm 的指标要根据业务语义定义。

\section{常见 benchmark 反模式}

\subsection{反模式 1：把打印放进 timed region}

错误：

\begin{lstlisting}[language=C++]
auto begin = Clock::now();
const auto result = kernel(input);
std::printf("%llu\n", static_cast<unsigned long long>(result));
auto end = Clock::now();
\end{lstlisting}

打印可能比 kernel 慢得多，也可能触发锁、缓冲、系统调用。

\subsection{反模式 2：每轮重新分配}

错误：

\begin{lstlisting}[language=C++]
for (int iter = 0; iter < repetitions; ++iter) {
    auto begin = Clock::now();
    std::vector<float> values(size);
    kernel(values);
    auto end = Clock::now();
}
\end{lstlisting}

这测了 allocator 和触页。除非目标就是测分配，否则分配应放在 setup。

\subsection{反模式 3：输入太小}

如果 kernel 只运行几十纳秒，计时器开销、函数调用、分支、循环控制都会淹没结果。解决方法：

\begin{itemize}
  \item 增大输入。
  \item 每次计时运行多个 batch。
  \item 单独测计时器 overhead。
  \item 使用硬件 counter 或专门 benchmark 框架。
\end{itemize}

\subsection{反模式 4：只测一次}

一次运行可能被中断、调度、缺页、降频影响。至少重复几十次，并保存原始样本。

\subsection{反模式 5：比较不同编译参数}

如果 A 用 \code{-O2}，B 用 \code{-O3 -march=native}，你比较的不是代码差异，而是代码加编译策略的组合。除非这正是实验目标，否则编译参数必须一致。

\subsection{反模式 6：没有审计汇编}

如果你没有看反汇编，就不知道测的是手写循环、自动向量化循环、库函数调用，还是被优化器删掉的空代码。

\section{报告模板}

每次实验至少提交：

\begin{lstlisting}
title:
    benchmark name and purpose

machine:
    CPU model
    memory if relevant
    OS and kernel
    governor / taskset / VM status if known

compiler:
    compiler version
    full flags

source:
    commit hash or file version
    kernel signature
    reference implementation

workload:
    input size
    data distribution
    cache state
    repetitions
    warmup count

measurement:
    timer or perf command
    raw data file
    summary statistics

binary audit:
    relevant objdump excerpt
    whether function was inlined/vectorized

correctness:
    validation method
    tolerance for floating point

conclusion:
    what the data supports
    what remains uncertain
\end{lstlisting}

\begin{keyidea}
高质量性能报告不是“我的版本更快”。它应该让别人能复现实验、审计二进制、理解误差，并判断结论的适用范围。
\end{keyidea}

\section{实验 15.1：修复一个会被优化器删除的 benchmark}

\begin{labbox}
目标：亲手观察 dead-code elimination 和防优化工具的作用。

创建 \filepath{labs/ch15_measurement/lab15_1_dce/bench_bad.cpp}：

\begin{lstlisting}[language=C++]
#include <cstdint>
#include <vector>

int main()
{
    constexpr std::uint64_t LOOP_LIMIT = 1000000000ULL;
    std::uint64_t sum = 0;
    for (std::uint64_t i = 0; i < LOOP_LIMIT; ++i) {
        sum += i;
    }
}
\end{lstlisting}

任务：

\begin{enumerate}
  \item 用 \code{-O0}、\code{-O2}、\code{-O3} 分别编译。
  \item 用 \code{objdump -drwC -Mintel} 检查循环是否存在。
  \item 加入返回值打印，再看汇编是否变成公式或仍保留循环。
  \item 把输入上限改成运行时参数，例如从 \code{argv} 读取，再观察。
  \item 加入 \code{do_not_optimize(sum)}，解释汇编变化。
\end{enumerate}

命令：

\begin{lstlisting}[language=bash]
clang++ -std=c++20 -O0 bench_bad.cpp -o bench_O0
clang++ -std=c++20 -O3 bench_bad.cpp -o bench_O3
objdump -drwC -Mintel bench_O3 > bench_O3.objdump.txt
rg -n "add|loop|main" bench_O3.objdump.txt
\end{lstlisting}

交付物：

\begin{enumerate}
  \item 三个优化级别的反汇编关键片段。
  \item 哪个版本保留循环，哪个版本删除或改写循环。
  \item 用 C++ 抽象机可观察行为解释为什么优化器有权这么做。
  \item 修复后的 benchmark。
\end{enumerate}
\end{labbox}

\section{实验 15.2：写一个最小可信 sum benchmark}

\begin{labbox}
目标：实现本章最小 harness，并形成第一份性能报告。

要求：

\begin{enumerate}
  \item 使用 \code{std::chrono::steady_clock}。
  \item setup、warmup、timed region、validation 分开。
  \item 输入规模至少包含：
\begin{lstlisting}
1024
16384
262144
1048576
16777216
\end{lstlisting}
  \item 每个规模 warmup 5 次，正式测量 31 次。
  \item 输出 CSV：
\begin{lstlisting}
n,min_ns,median_ns,max_ns,bytes,bandwidth_GBps
\end{lstlisting}
  \item 用 \code{objdump} 审计核心循环。
\end{enumerate}

交付物：

\begin{enumerate}
  \item 源码和编译命令。
  \item CSV 原始输出。
  \item 一张结果表。
  \item 汇编核心循环标注。
  \item 500 字报告：随着 \code{n} 增大，瓶颈可能如何变化。
\end{enumerate}
\end{labbox}

\section{实验 15.3：比较 \code{steady_clock} 与 TSC}

\begin{labbox}
目标：理解时间源差异，而不是盲目相信某一个计时器。

任务：

\begin{enumerate}
  \item 对同一个 \code{sum_u64} kernel，同时记录 \code{steady_clock} 纳秒和 TSC delta。
  \item 用 \instruction{rdtsc} 裸读一次。
  \item 用 \instruction{lfence} 包围 \instruction{rdtsc} 再读一次。
  \item 每种方式重复 100 次，保存原始数据。
  \item 检查线程是否迁核；如果会迁核，用 \code{taskset} 绑定。
\end{enumerate}

交付物：

\begin{enumerate}
  \item 计时代码。
  \item 原始数据。
  \item min、median、max 表。
  \item 解释 TSC delta 是否稳定。
  \item 说明为什么 TSC delta 不应无条件等同于真实 core cycles。
\end{enumerate}
\end{labbox}

\section{实验 15.4：用 \code{perf stat} 建立硬件证据}

\begin{labbox}
目标：从单纯时间测量进入硬件 counter 观察。

对实验 15.2 的 benchmark 运行：

\begin{lstlisting}[language=bash]
perf stat -r 10 \
    -e cycles,instructions,branches,branch-misses,cache-references,cache-misses \
    ./bench_sum
\end{lstlisting}

任务：

\begin{enumerate}
  \item 记录 \code{cycles}、\code{instructions}、IPC。
  \item 记录 cache miss 相关事件。
  \item 如果权限不足，记录完整错误信息和 \code{perf_event_paranoid} 值。
  \item 改变输入规模，观察事件是否随 \code{n} 变化。
  \item 把 PMU 结果和时间结果放在同一份报告中解释。
\end{enumerate}

交付物：

\begin{enumerate}
  \item \code{perf stat} 输出。
  \item 事件含义解释。
  \item 输入规模与 counter 变化表。
  \item 至少一个你不能确定的点，并说明还需要什么实验。
\end{enumerate}
\end{labbox}

\section{实验 15.5：cache 状态实验}

\begin{labbox}
目标：观察 warm-cache 与 cold-ish-cache 对结果的影响。

设计两个版本：

\begin{enumerate}
  \item warm-cache：同一数组连续重复求和。
  \item cold-ish-cache：每次计时前扫描一个远大于 LLC 的 buffer，尽量污染 cache。
\end{enumerate}

示例 cache 污染函数：

\begin{lstlisting}[language=C++]
void touch_buffer(std::span<std::uint64_t> buffer)
{
    constexpr std::uint64_t TOUCH_INCREMENT = 1;
    for (std::uint64_t& value : buffer) {
        value += TOUCH_INCREMENT;
    }
}
\end{lstlisting}

任务：

\begin{enumerate}
  \item 选择至少 3 个输入规模。
  \item 每个规模分别测 warm-cache 和 cold-ish-cache。
  \item 报告 median ns 和 GB/s。
  \item 解释为什么 cold-ish-cache 只是近似，不是严格清空所有 cache。
  \item 用 \code{perf stat} 观察 cache miss 是否变化。
\end{enumerate}

交付物：

\begin{enumerate}
  \item 源码。
  \item 结果表。
  \item cache 状态解释。
  \item 对后续内存层次章节提出 3 个问题。
\end{enumerate}
\end{labbox}

\section{本章作业}

\begin{exercisebox}
\subsection*{作业 1：benchmark 批判性阅读}

找一个你过去写过的性能测试，或者网上任意一段短 benchmark。写一份审计报告，回答：

\begin{enumerate}
  \item 被测对象是否明确？
  \item setup 和 timed region 是否分离？
  \item 是否有 correctness reference？
  \item 是否可能被 DCE、constant folding 或 loop invariant hoisting 影响？
  \item 是否重复测量？
  \item 是否报告统计口径？
  \item 是否记录环境？
  \item 是否审计汇编？
\end{enumerate}

要求至少指出 5 个问题，并给出修复方案。

\subsection*{作业 2：最小 harness 代码评审}

阅读本章最小 C++20 harness，写一份代码评审：

\begin{itemize}
  \item 哪些设计是为了防止编译器优化掉被测代码？
  \item 哪些设计是为了避免把 setup 混进计时区域？
  \item 哪些常量应该暴露成命令行参数？
  \item \code{median_ns} 当前实现有什么限制？
  \item 如果要支持多个 kernel，应该如何拆分函数和数据结构？
\end{itemize}

要求给出不少于 800 字的工程评审，不要只复述代码。

\subsection*{作业 3：时间源误差分析}

对一个空 timed region 测量计时器开销：

\begin{lstlisting}
begin = now()
end = now()
\end{lstlisting}

再对一个很短 kernel 和一个较长 kernel 分别计时。要求：

\begin{enumerate}
  \item 报告空计时开销分布。
  \item 解释为什么短 kernel 更容易被计时器开销污染。
  \item 设计 batch 策略降低相对误差。
  \item 说明 batch 策略可能引入什么新问题。
\end{enumerate}

\subsection*{作业 4：汇编审计报告}

选择实验 15.2 中一个输入规模，对可执行文件生成反汇编报告：

\begin{lstlisting}[language=bash]
objdump -drwC -Mintel ./bench_sum > bench_sum.objdump.txt
\end{lstlisting}

报告必须包含：

\begin{enumerate}
  \item \code{sum_u64} 是否被内联。
  \item 核心循环反汇编。
  \item 是否出现 SIMD 指令。
  \item 每次循环处理多少元素。
  \item 循环退出条件。
  \item 源码中的 \code{do_not_optimize} 在汇编层是否生成真实指令。
\end{enumerate}

\subsection*{作业 5：性能实验报告}

提交一份完整性能实验报告，主题是 \code{sum_u64}、\code{copy_u64} 或 \code{dot_f32} 三选一。报告必须包含：

\begin{itemize}
  \item 环境记录。
  \item 编译命令。
  \item workload 定义。
  \item correctness reference。
  \item 原始测量数据。
  \item min、median、p95、max。
  \item 至少一段反汇编。
  \item 至少一次 \code{perf stat} 尝试。
  \item 结论和限制。
\end{itemize}

\subsection*{评分标准}

本章作业按下面标准评价：

\begin{enumerate}
  \item 是否能清楚划分 setup、warmup、timed region 和 validation。
  \item 是否能解释编译器删除、移动或常量折叠 benchmark 的原因。
  \item 是否能用 reference 证明结果正确。
  \item 是否能保存和解释重复测量数据。
  \item 是否能选择合适指标，例如 ns/op、GB/s、GFLOP/s。
  \item 是否能用 \code{objdump} 和 \code{perf stat} 支撑结论。
  \item 是否能诚实说明实验限制，而不是过度推断。
\end{enumerate}
\end{exercisebox}

\section{本章小结}

本章建立了性能测量的第一层地基：

\begin{lstlisting}
correctness:
    reference first, performance second

benchmark structure:
    setup -> warmup -> timed region -> validation -> report

compiler safety:
    avoid DCE, constant folding, unintended motion

time source:
    steady_clock for basic timing
    TSC and PMU require stronger interpretation

noise:
    OS scheduling, CPU frequency, cache state, migration, background load

evidence:
    raw samples, statistics, objdump, perf stat, environment record
\end{lstlisting}

从现在开始，你不应该再接受“这段代码跑了几毫秒”这种孤立数字。性能工程的基本素养是：每个数字都必须带着上下文、方法、证据和限制。下一章会继续把这些原则工程化，系统设计 benchmark 输入、热路径、CSV 输出、参数扫描和报告格式。
